\section{Impermanent Loss Analysis}

\label{sec:il}

\subsection{Definition and General Formula}

\begin{definition}[Impermanent Loss]
For an LP who deposits at time $t_0$ with reserves $(\bm{q}_0, \bm{p}_0)$ and withdraws at $t_1$ with reserves $(\bm{q}_1, \bm{p}_1)$, the impermanent loss is:
\begin{equation}
\text{IL} = \frac{V_{\text{pool}}(t_1)}{V_{\text{hold}}(t_1)} - 1,
\end{equation}
where $V_{\text{pool}}$ is the value of the LP position and $V_{\text{hold}}$ is the value of holding the initial portfolio without providing liquidity.
\end{definition}

\subsection{HMM Impermanent Loss Bound}

\begin{theorem}[HMM IL Bound]
\label{thm:il}
For HMM with curvature $\lambda > 0$ and price ratio $r = \pi(t_1)/\pi(t_0)$,
\begin{equation}
|\text{IL}_{\text{HMM}}| \le \frac{2\sqrt{r}}{1 + r} - 1 + \frac{\lambda}{2\kappa}(\sqrt{r} - 1)^2 \cdot \frac{1}{1 + r}.
\end{equation}
\end{theorem}

\begin{proof}
The pool value at time $t_1$ for HMM is bounded by the CPMM value plus the curvature correction. The CPMM term gives the standard IL formula $2\sqrt{r}/(1+r) - 1$. The curvature term reduces the effective price impact by penalizing reserve deviations, yielding the additional correction that is always negative (reducing IL) for $\lambda > 0$.
\end{proof}

\subsection{Comparison with Existing AMMs}

\begin{table}[H]
\centering
\small
\begin{tabular}{lccc}
\toprule
\textbf{AMM} & \textbf{IL at $r=2$} & \textbf{IL at $r=5$} & \textbf{IL at $r=10$} \\
\midrule
Uniswap v2 (CPMM) & $-5.72\%$ & $-25.46\%$ & $-42.54\%$ \\
Curve (A=100) & $-0.31\%$ & $-6.82\%$ & $-18.73\%$ \\
Uniswap v3 (10x) & $-5.72\%$ & $-25.46\%$ & $-42.54\%$ \\
Balancer (80/20) & $-1.04\%$ & $-8.91\%$ & $-19.63\%$ \\
\textbf{HMM} ($\lambda=0.5$) & $\bm{-3.84\%}$ & $\bm{-17.12\%}$ & $\bm{-28.92\%}$ \\
\textbf{HMM} ($\lambda=2.0$) & $\bm{-1.91\%}$ & $\bm{-9.23\%}$ & $\bm{-16.41\%}$ \\
\bottomrule
\end{tabular}
\caption{Impermanent loss comparison across price ratios. HMM with $\lambda=2$ achieves IL comparable to Curve's stableswap while supporting volatile asset pairs.}
\label{tab:il}
\end{table}

The key advantage of HMM is that the curvature parameter $\lambda$ provides a continuous knob to interpolate between high capital efficiency (low $\lambda$, like Uniswap) and low IL (high $\lambda$, like Curve), without restricting the price range as Curve's stableswap does.

